\label{sec:intro}

Large language models (LLMs) are increasingly used for software engineering tasks such as code generation, code review, and vulnerability detection. In security-related tasks, an LLM is often asked to inspect a code snippet and decide whether it contains a known type of vulnerability, such as a Common Weakness Enumeration (CWE) issue. This setting is useful because it can help developers find bugs earlier, but it is also risky because an incorrect prediction can have practical consequences. A false positive may waste developer time by reporting safe code as vulnerable, while a false negative may miss a real security weakness. Recent benchmarks such as CASTLE evaluate LLMs, static analyzers, and formal verification tools on CWE detection using hand-crafted C micro-benchmarks, showing both the promise and limitations of LLM-based vulnerability detection~\cite{castle,castle_github}. CASTLE contains 250 hand-crafted C micro-benchmark programs covering 25 common CWE categories, which makes it suitable for controlled testing of vulnerability detection behavior~\cite{castle_github}.

A standard LLM baseline treats each test case independently. The model receives a code sample, produces a prediction, and then the interaction ends. Even if the model makes a mistake, that mistake is not preserved as part of the next inference step. This creates a gap between human debugging behavior and standard LLM inference. A human developer can remember previous mistakes and use them as warnings when seeing similar code later, but a normal LLM pipeline has no persistent memory of its own past failures. As a result, the same model may repeatedly over-report safe code patterns or repeatedly miss similar vulnerable patterns. This repeated-error problem is especially important in CWE detection because many vulnerability categories share recurring code structures.

In this paper, we introduce \system, a database-backed mistake-memory framework for LLM vulnerability detection. The key idea is to store incorrect model predictions in a persistent mistake database and retrieve similar past mistakes when analyzing a new code sample. Instead of fine-tuning the model or changing the model architecture, \system changes the inference pipeline. First, a baseline LLM is evaluated on labeled examples. Second, the incorrect predictions are inserted into a SQLite database with structured metadata, including the sample identifier, source code, true label, predicted label, true CWE, predicted CWE, raw model output, and error type. Third, during future inference, \system retrieves similar stored mistakes using TF-IDF similarity and injects them into the prompt as cautionary examples. SQLite is appropriate for this prototype because it is lightweight, self-contained, and does not require a separate database server~\cite{sqlite}.

This design is different from a standard LLM baseline because it adds a persistent feedback loop. The baseline only reasons from the current prompt and the model's internal parameters. \system additionally reasons from a model-specific external memory of previous failures. This external memory makes the system more inspectable: each ErrorLock prediction can be traced to the retrieved mistake identifiers and similarity scores used in the augmented prompt. This is useful because the benchmarks used in our project provide labels, CWE categories, and vulnerable-line information, but they do not provide full execution traces or human reasoning traces for each sample. Therefore, \system creates its own lightweight trace of model behavior through stored prediction records and retrieval metadata.

We evaluate \system using Qwen2.5-Coder-7B-Instruct on two CWE-based benchmarks. CASTLE-C250 is used as the main controlled benchmark. To evaluate benchmark scope and generalizability beyond the original mistake source, we also construct a sampled Juliet C/C++ Test Suite subset whose CWE categories overlap with CASTLE. Juliet is a larger NIST test suite organized under 118 different CWEs~\cite{juliet_nist}. In our original experiments, the ErrorLock database is built from CASTLE mistakes and then tested both on CASTLE and on the external Juliet subset. This setup allows us to study whether mistake memory only helps within the original benchmark or can also transfer to a different CWE benchmark. We additionally include a supplemental full-C250 experiment using OpenRouter's \texttt{openai/gpt-oss-120b:free} model and a compact RAG memory variant to test whether stronger retrieval and a stronger model improve the held-out CASTLE result.

This project currently focuses on CWE-based vulnerability detection rather than general software engineering tasks. However, \system is not tied to a specific benchmark format. Its core mechanism is to store model mistakes, retrieve similar past failures, and use them as prompt context for future predictions. In this report, we evaluate this idea on CASTLE as the main controlled benchmark and Juliet as an external benchmark. The benchmarks provide labels and CWE metadata, but not full reasoning traces; therefore, \system records its own lightweight traces through stored prediction records, retrieved mistake identifiers, and similarity scores.

Specifically, this paper makes the following contributions:
\begin{itemize}
    \item We propose \system, a database-backed mistake-memory framework that stores incorrect LLM vulnerability predictions and retrieves similar past mistakes during future inference.
    \item We implement a prototype using SQLite as the backend mistake database, TF-IDF similarity for retrieval, and retrieval-augmented prompting for prediction correction.
    \item We evaluate the prototype on CASTLE-C250 and a sampled Juliet C/C++ subset. Our results show mixed behavior on CASTLE, where \system reduces false positives but can increase false negatives, and a stronger improvement on Juliet, where accuracy improves from 68.67\% to 82.40\%.
    \item We add a supplemental full-C250 compact-RAG experiment with \texttt{openai/gpt-oss-120b:free}, where \system improves held-out F1 from 0.616 to 0.709 and CWE correctness from 0.433 to 0.520.
\end{itemize}

The remainder of this paper is organized as follows. Section~\ref{sec:related} reviews related work on CWE benchmarks, LLM-based vulnerability detection, and retrieval-augmented methods. Section~\ref{sec:overview} describes the design and implementation of \system. Section~\ref{sec:eval} presents the experimental setup and results. Finally, Section~\ref{sec:conclude} concludes the paper and discusses limitations and future work.
